
% SGA 3，报告 VII_B；英文重建本第一轮的中文译本。
% 范围： 节 2.6.3 全文 依据文献局部页 62,
% 止于下节之前： 节 2.7.
% 合并阅读本页： 575. 重版印刷页： 564.
% 范围内原版边码： 529.
% 范围内编者注： 111.
% 依据： Polo--Gille Exp7B-13oct24.pdf,
% SHA-256 D303690E0195654BBE7C3DAA5D89EA23E24CD709A3575217234AC19AE8836694.
% 参照了 Jacob Reinhold 的 CC BY 4.0 Markdown，作为对照来源。
% 原英文 Loop 1 阶段说明：交换图为第一轮直接从 PDF 截取的原图。（当前中文工程已采用原生图。）
\subsubsection*{2.6.3}
\phantomsection
\label{sga3:VIIB:2.6.3}
若 \(\mathbf L\) 是 \(\mathbf O_k\) 上的平坦 Lie 代数，则形式群 \(\mathscr G(\mathbf L)\) 可由一个泛性质刻画。\footnote{%
\label{sga3:VIIB:nde:111}编者修改了以下内容，以利用 \SGARefLink{sga3:VIIB:2.6.2:prop}{2.6.2} 中的补充。}
事实上，每个到形式群 \(G\) 的态射
\[
  \phi:\mathscr G(\mathbf L)\longrightarrow G
\]
诱导一个态射
\[
  \operatorname{Lie}(\phi):
  \operatorname{Lie}\bigl(\mathscr G(\mathbf L)\bigr)
  \longrightarrow
  \operatorname{Lie}(G).
\]
将其与
\[
  \tau_{\mathbf L}:
  \mathbf L\longrightarrow
  \operatorname{Lie}\bigl(\mathscr G(\mathbf L)\bigr)
\]
\textup{（参见\ \SGARefLink{sga3:VIIB:2.6.2:prop}{2.6.2}）}复合，得到态射
\[
  \phi':\mathbf L\longrightarrow\operatorname{Lie}(G).
\]
于是有：
\medskip
\noindent\textbf{命题。}
\phantomsection
\label{sga3:VIIB:2.6.3:prop}
{\itshape 设 \(G\) 是形式 \(k\)-群，\(\mathbf L\) 是 \(\mathbf O_k\) 上的平坦 Lie 代数。上面定义的映射 \(\phi\mapsto\phi'\) 是双射
\[
  \operatorname{Hom}_{\operatorname{Grf}/k}
    \bigl(\mathscr G(\mathbf L),G\bigr)
  \xrightarrow{\ \sim\ }
  \operatorname{Hom}_{\operatorname{Lie}}
    \bigl(\mathbf L,\operatorname{Lie}(G)\bigr),
\]
其中右边是从 \(\mathbf L\) 到 \(\operatorname{Lie}(G)\) 的 \(\mathbf O_k\) 上 Lie 代数态射的集合。}
\textit{证明。}
可立即归约到 \(k\) 是 Artin 环的情形。置 \(L=\mathbf L(k)\)。由 \SGARefLink{sga3:VIIB:2.3.1}{2.3.1}，
\[
  \operatorname{Hom}_{\operatorname{Grf}/k}
    \bigl(\mathscr G(\mathbf L),G\bigr)
\]
与使下图交换的含幺代数态射
\[
  h:U(L)\longrightarrow\mathbf H(G)
\]
的集合存在双射：
\begin{center}
  \phantomsection
  \label{sga3:VIIB:2.6.3:universal-property-diagram}
  % SGA3-VII-DIAG-119
  % SGA3-VII-DIAG-119
% Authority: Exp7B-13oct24.pdf, local p. 62.
% Crop: figures/Exp7B_localp062_universal_property_commutative_diagram.png
% Locator: tex/components/53_expose_VIIB_section263_universal_property_en.tex, Loop 1 line 77.
\begin{tikzcd}[row sep=1.2em,column sep=4.2em]
U(L) \arrow[r,"h"] \arrow[dd,"\Delta_{U(L)}"']
  & \mathbf H(G) \arrow[dr,"\delta_G"] & {} \\
{} & {} & \mathbf H(G\times G) \\
U(L)\otimes U(L) \arrow[r,"h\otimes h"]
  & \mathbf H(G)\otimes\mathbf H(G)
      \arrow[ur,"\varphi_G"']
  & {.}
\end{tikzcd}

\end{center}
现在，\(h\) 由其在 \(L\) 上的限制决定，这个限制是从 \(L\) 到 \(\mathbf H(G)\) 的底 Lie 代数的态射。该图交换，恰好是说 \(h\) 将 \(L\) 映入 \(\mathbf H(G)\) 中满足
\[
  \delta_G(x)
  =
  \varphi_G(x\otimes1+1\otimes x).
\]
的元素 \(x\) 构成的子集。这个子集恰好是 \(\operatorname{Lie}(G)\)；参见\ \SGARefLink{sga3:VIIB:2.6:item:c}{2.6(c)}。命题得证。
